%----------------------------------------------------------------------------------------
%	VS CODE DARK MODERN THEME FOR LATEX LISTINGS
%----------------------------------------------------------------------------------------

% --- COLORES OFICIALES DEL TEMA VS CODE "DARK MODERN" ---
\definecolor{vscBack}{RGB}{30,30,30}          % Fondo oscuro del editor (#1e1e1e)
\definecolor{vscText}{RGB}{204,204,204}        % Texto normal (#cccccc)
\definecolor{vscKeyword}{RGB}{86,156,214}      % Palabras clave estándar - azul (#569cd6)
\definecolor{vscControl}{RGB}{197,134,192}     % Control de flujo (if, return) - rosa/morado (#c586c0)
\definecolor{vscType}{RGB}{78,201,176}         % Clases, estructuras y tipos - turquesa (#4ec9b0)
\definecolor{vscString}{RGB}{206,145,120}       % Cadenas/Strings - salmón/naranja (#ce9178)
\definecolor{vscComment}{RGB}{106,153,85}       % Comentarios - verde apagado (#6a9955)
\definecolor{vscNumber}{RGB}{181,206,168}       % Dígitos y números - verde claro (#b5cea8)
\definecolor{vscFunction}{RGB}{220,220,170}     % Métodos y nombres de función - amarillo (#dcdcaa)
\definecolor{vscGutter}{RGB}{133,133,133}       % Color para números de línea (#858585)
\definecolor{vscFrame}{RGB}{60,60,60}           % Bordes/Marco del código (#3c3c3c)

% --- CONFIGURACIÓN DE FUENTES (Sugerencia) ---
% Si compilas con pdfLaTeX:
% Puedes instalar/usar el clon libre de Consolas llamado 'inconsolata':
% \usepackage[varqu]{zi4} 
%
% Si compilas con XeLaTeX o LuaLaTeX (Recomendado para fuentes del sistema):
% Puedes usar directamente la tipografía Consolas de tu máquina:
% \usepackage{fontspec}
% \setmonofont{Consolas}

% --- DEFINICIÓN DE LENGUAJE C# ESTILO VS CODE ---
\lstdefinelanguage{CSharpVSCode}{
	% Clase 1: Palabras clave estándar (Azul)
	morekeywords=[1]{
		using, namespace, class, public, private, protected, internal,
		readonly, static, void, string, int, bool, object, out, var,
		new, null, byte, short, long, float, double, decimal, char,
		struct, enum, delegate, event, interface, get, set, value,
		this, override, virtual, abstract, partial, false, true
	},
	% Clase 2: Control de flujo y modificadores (Rosa/Morado)
	morekeywords=[2]{
		if, else, while, for, foreach, try, catch, finally, throw, return,
		break, continue, lock, await, async, in, switch, case, default
	},
	% Clase 3: Clases y Tipos específicos (Turquesa)
	morekeywords=[3]{
		SerialService, SerialHub, IHubContext, ILogger, SerialPort, Thread,
		Exception, TimeoutException, Task, Path, AppDomain, File,
		JsonSerializer, JsonSerializerOptions, SettingsData, Array,
		StringComparison, Action, Console, TaskRun
	},
	sensitive=true,
	morecomment=[l]{//},
	morecomment=[s]{/*}{*/},
	morestring=[b]",
	morestring=[b]'
}

% --- ESTILO DE LISTING PRINCIPAL (VSCodeDark) ---
\lstdefinestyle{VSCode}{
	language=CSharpVSCode,
	backgroundcolor=\color{vscBack},
	basicstyle=\color{vscText}\ttfamily\small,
	keywordstyle=[1]\color{vscKeyword}\bfseries,
	keywordstyle=[2]\color{vscControl}\bfseries,
	keywordstyle=[3]\color{vscType},
	stringstyle=\color{vscString},
	commentstyle=\color{vscComment}\itshape,
	breakatwhitespace=false,
	breaklines=true,
	captionpos=b,
	keepspaces=true,
	numbers=left,
	numbersep=10pt,
	numberstyle=\tiny\color{vscGutter},
	showspaces=false,
	showstringspaces=false,
	showtabs=false,
	tabsize=4,
	frame=single,
	rulecolor=\color{vscFrame},
	frameround=tttt,
	columns=flexible,
	% Regla de reemplazo para pintar números enlistings
	literate=*{\%}{{{\color{vscControl}\%}}}{1}
		{0}{{{\color{vscNumber}0}}}{1}
		{1}{{{\color{vscNumber}1}}}{1}
		{2}{{{\color{vscNumber}2}}}{1}
		{3}{{{\color{vscNumber}3}}}{1}
		{4}{{{\color{vscNumber}4}}}{1}
		{5}{{{\color{vscNumber}5}}}{1}
		{6}{{{\color{vscNumber}6}}}{1}
		{7}{{{\color{vscNumber}7}}}{1}
		{8}{{{\color{vscNumber}8}}}{1}
		{9}{{{\color{vscNumber}9}}}{1}
}
